(1) 初态总能量 \(E_i=(\gamma+1)mc^2\)，动量 \(p_i=\gamma m v\)，因此末态复合物速度由 \(V=c^2p_i/E_i\) 给出
\[
V=\frac{\gamma}{\gamma+1}v,\qquad
\frac Vc=\frac{\gamma\beta}{\gamma+1}
=\sqrt{\frac{\gamma-1}{\gamma+1}}.
\]
由 \(1-V^2/c^2=2/(\gamma+1)\) 得 \(\Gamma=\sqrt{(\gamma+1)/2}\)，其中 \(\Gamma\) 是复合物的实验室洛伦兹因子。

(2) 末态静质量由不变量
\[
M^2c^4=E_i^2-c^2p_i^2
=m^2c^4[(\gamma+1)^2-\gamma^2\beta^2]
=2m^2c^4(\gamma+1).
\]
所以
\[
M=m\sqrt{2(\gamma+1)}>2m
\]
因为 \(\gamma>1\)。复合物内部能使静质量增加，不能仍把其惯性质量强行视作 \(2m\)。

(3) 以两粒子各自的静能之和为基准，复合物内部能增量
\[
\Delta U=(M-2m)c^2=mc^2[\sqrt{2(\gamma+1)}-2]>0.
\]
末态总能量为 \(\Gamma Mc^2=\sqrt{(\gamma+1)/2}\cdot m\sqrt{2(\gamma+1)}c^2=(\gamma+1)mc^2\)，总动量为 \(\Gamma M V=E_iV/c^2=\gamma mv\)，分别与初态一致。实验室宏观动能的减少等于复合物静能增加，但不可把非弹性碰撞中机械动能守恒作为前提。

(4) 当 \(\beta\ll1\)，\(\gamma=1+\beta^2/2+O(\beta^4)\)，故
\[
V=\frac v2+O(v^3/c^2),\qquad
M=2m\left(1+\frac{\beta^2}{8}+O(\beta^4)\right),
\quad
\Delta U=\frac{mv^2}{4}+O(mv^4/c^2).
\]
经典两等质量粘连碰撞末态速度 \(v/2\)、动能损失 \(mv^2/4\)，与极限完全吻合。
